\begin{abstract}
Urban new energy logistics scheduling is more complex than finding the shortest route. In practice, vehicles have to make decisions while tasks keep arriving, battery and load capacity are limited, charging takes time, stations may have queues, and unfinished tasks may miss their deadlines. To study this process, we built a multi-vehicle scheduling simulator and used it to compare different dispatching strategies.

The road network is represented as a weighted graph. Dijkstra and A* are used for shortest-path and feasibility queries, while RRT is kept as an extensible path-planning module. For scheduling, we implement Nearest, EDF, Heaviest, a Q-learning based rule selector, and an offline MILP reference for small-scale full-information cases. The simulator is connected with a backend service and a visualization frontend, so the scheduling process can be executed, observed, and compared.

Experiments on controlled random graph scenarios show that Nearest is a strong fixed baseline, while Q-learning can learn useful rule-selection behavior but is less stable in medium-scale settings. As an extension, we also build a processed Guangzhou Panyu OSM road-network scenario to test the same simulator on a real urban map. The focus is a runnable, visible, and comparable scheduling platform rather than a single algorithm.
\end{abstract}
